\section{Universal nondegeneracy}\label{sec:nondegeneracy}

\Cref{sec:local-reduction} reduced the optimization problem
\eqref{eq:graphon-variational-problem} on the symmetry-breaking side to a one-variable
problem in the edge density $\eps=e(W)$.  We now study the behaviour of the
reduced objective $\mathcal F_{p,r}(\eps)$ near the boundary value in terms of
the edge deficit
\[
  \delta=r-\eps.
\]
Note that by \Cref{lem:edge-density-deficit}, $\delta>0$ for every optimizer on the
symmetry-breaking side.
At the boundary value $p_0=\pc(r)$, define the boundary excess
\begin{equation}\label{eq:boundary-excess}
  G_r(\delta)=\mathcal F_{p_0,r}(r-\delta)-J_{p_0}(r),
  \qquad 0<\delta<\rho_0,
\end{equation}
with $\rho_0$ as in \Cref{lem:fixed-density-bipodality}; that lemma ensures
that $\mathcal F_{p_0,r}(r-\delta)$ is finite throughout this range.  The value
at $\delta=0$ belongs instead to the analytic extension $G(r,\delta)$
constructed in \Cref{thm:positive-second-variation}; it is not part of the domain of
$G_r$.
Thus, $G_r(\delta)$ is the excess $I_{p_0}$-value, relative to the constant
graphon, incurred at the boundary by forcing edge density $r-\delta$ while
keeping the $H$-density at the same level $r^m$.  From now on, a subscript $K$
in asymptotic notation means
that the implicit constants may depend on $K$, but are uniform for $r\in K$.
The goal is to prove that this boundary excess has a
nondegenerate quadratic minimum at $\delta=0$:
\[
  G_r(\delta)=\frac12A_H(r)\delta^2+O_K(\delta^3)
  \quad\text{as }\delta\downarrow0,
  \qquad A_H(r)>0.
\]

\subsection{A scalar quadratic lower bound}

Recall that $K$ is a compact subinterval of $(0,1)\setminus\{r_*\}$.  Fix
$r\in K$, and recall that $\sm(r)$ is the second contact density associated
with the boundary point $(\pc(r),r)$.  Let $\ell_r$ be the supporting line to
the graph of $x\mapsto J_{\pc(r)}(x^{1/d})$ at $x=r^d$, namely
\[
  \ell_r(x)=J_{\pc(r)}(r)
  +\frac{J_{\pc(r)}'(r)}{d r^{d-1}}(x-r^d).
\]
Define
\begin{equation}\label{eq:supporting-cost-gap}
  g_r(u)=J_{\pc(r)}(u)-\ell_r(u^d),
  \qquad 0\le u\le1.
\end{equation}
Since $\ell_r$ is a supporting line, $g_r\ge0$.  We claim that there is
$\gamma>0$, uniform for $r\in K$, such that
\begin{equation}\label{eq:supporting-gap-quadratic-bound}
  g_r(u)\ge \gamma\,\dist(u,\{r,\sm(r)\})^2.
\end{equation}
To prove the claim, we first use the quadratic separation condition
\eqref{eq:contact-quadratic-separation}, which gives
\[
  g_r(u)\ge \gamma_0\,\dist(u^d,\{r^d,\sm(r)^d\})^2
\]
with $\gamma_0>0$ uniform for $r\in K$.
The two contact densities stay in a compact subset of $(0,1)$, and thus there
exists $\eta_K>0$ such that $v\ge \eta_K$ for all $r\in K$ and
$v\in\{r,\sm(r)\}$.
For any $u\in[0,1]$,
\[
  |u^d-v^d|
  =
  |u-v|\,(u^{d-1}+u^{d-2}v+\cdots+v^{d-1})
  \ge
  \eta_K^{d-1}|u-v|.
\]
Taking the closer of the two contact densities gives
\[
  \dist(u^d,\{r^d,\sm(r)^d\})
  \ge
  \eta_K^{d-1}\dist(u,\{r,\sm(r)\}).
\]
Thus, after decreasing the constant, we obtain \eqref{eq:supporting-gap-quadratic-bound}.  We
also note that, since $K$ is compact and avoids $r_*$, continuity of $\sm$ and
the fact that $\sm(r)\ne r$ imply that $|\sm(r)-r|$ is bounded away from zero
on $K$.

\begin{lemma}[Scalar quadratic lower bound]\label{lem:scalar-quadratic-bound}
There exist constants $C>0$ and $\delta_0>0$, uniform for $r\in K$, such that
the following holds.  If $0<\delta<\delta_0$ and $X$ is a random variable with
values in $[0,1]$ satisfying
$\EE X=r-\delta$ and
$\EE X^d\ge r^d$,
then
\begin{equation}\label{eq:scalar-quadratic-bound}
  \EE J_{\pc(r)}(X)\ge J_{\pc(r)}(r)+C\delta^2.
\end{equation}
\end{lemma}

\begin{proof}
The slope of $\ell_r$ is positive, because $\pc(r)<r$.  Hence
\[
  \EE\ell_r(X^d)
  =   \ell_r(\EE X^d)
  \ge \ell_r(r^d)
  =   J_{\pc(r)}(r).
\]
Using \eqref{eq:supporting-cost-gap},
\[
  \EE g_r(X)
  =   \bigl(\EE J_{\pc(r)}(X)-\EE \ell_r(X^d)\bigr)
  \le \bigl(\EE J_{\pc(r)}(X)-J_{\pc(r)}(r)\bigr).
\]
It remains to show $\EE g_r(X)\ge C\delta^2$.  The estimate
\eqref{eq:supporting-gap-quadratic-bound} will force any counterexample to concentrate, up to
an $o(\delta)$ error in first moment, near the contact set $\{r,\sm(r)\}$.
If the second contact density lies above $r$, this is already incompatible
with the mean constraint $\EE X=r-\delta$.  If it lies below $r$, the mean
constraint determines to first order the mass near $\sm(r)$, and the resulting
loss in the $d$th moment contradicts $\EE X^d\ge r^d$.  We now make this
argument precise.

Suppose the bound fails.  Then, there exist $r_n\in K$, $\delta_n\downarrow0$,
and random variables $X_n$ with values in $[0,1]$ such that
\[
  \EE X_n=r_n-\delta_n,
  \qquad
  \EE X_n^d\ge r_n^d,
  \qquad
  \EE g_{r_n}(X_n)=o(\delta_n^2).
\]
Passing to a subsequence, assume $r_n\to\bar r$ and
$\sm(r_n)\to\bar s\ne\bar r$.  (The limit $\bar r\in K$ is not related to
the exceptional density $r_*$, which $K$ avoids.)  The separation
$\bar s\ne\bar r$ follows from the uniform separation noted above.  Let $Y_n$ be a measurable nearest-point
projection of $X_n$ onto the two-point set $\{r_n,\sm(r_n)\}$; that is, $Y_n$
is whichever of $r_n$ and $\sm(r_n)$ is closer to $X_n$, with ties chosen
arbitrarily.  By
\eqref{eq:supporting-gap-quadratic-bound},
\[
  \EE |X_n-Y_n|^2=o(\delta_n^2).
\]
By Jensen's inequality,
\[
  \bigl(\EE|X_n-Y_n|\bigr)^2
  \le \EE|X_n-Y_n|^2
  =o(\delta_n^2),
\]
and hence $\EE|X_n-Y_n|=o(\delta_n)$.
Thus, the projection changes the first moment only by $o(\delta_n)$.

Write
\[
  c_n=\mathbb P(Y_n=\sm(r_n)).
\]
Since $\EE X_n = \EE Y_n + \EE(X_n-Y_n)$ and
$|\EE(X_n-Y_n)| \le \EE|X_n-Y_n|=o(\delta_n)$, we have
\begin{equation}\label{eq:contact-projected-mean}
  r_n-\delta_n
  = \EE X_n
  = \EE Y_n+\EE(X_n-Y_n)
  =r_n+c_n(\sm(r_n)-r_n)+o(\delta_n).
\end{equation}
This equation immediately shows that $\bar s>\bar r$ is impossible.  To see
why, suppose that $\bar s>\bar r$.  Then $\sm(r_n)-r_n$ is bounded below by a
positive constant for all
large $n$.  Writing the error term in
\eqref{eq:contact-projected-mean} as $e_n=o(\delta_n)$, the right-hand side is at
least $r_n+e_n$.  For all large $n$, $e_n\ge -\delta_n/2$, and hence the
right-hand side is at least $r_n-\delta_n/2$, contradicting the left-hand side
$r_n-\delta_n$.

We now show that $\bar s<\bar r$ is also impossible.  If $\bar s<\bar r$, then the same
compactness and uniform separation give constants $\kappa,\eta,\eta'>0$ such
that, for all large $n$,
\[
  r_n-\sm(r_n)\ge \kappa,
  \qquad
  \eta \ge r_n^d-\sm(r_n)^d\ge \eta'.
\]
In particular, division by $r_n-\sm(r_n)$ is harmless, and
\eqref{eq:contact-projected-mean} implies
\[
  c_n=\frac{\delta_n}{r_n-\sm(r_n)}+o(\delta_n).
\]
Since $u\mapsto u^d$ is Lipschitz on $[0,1]$, with Lipschitz constant at most
$d$, the already proved estimate $\EE |X_n-Y_n|=o(\delta_n)$ gives
\[
  \left|\EE X_n^d-\EE Y_n^d\right|
  \le \EE |X_n^d-Y_n^d|
  \le d\,\EE |X_n-Y_n|
  =o(\delta_n).
\]
Therefore
\begin{align*}
  \EE X_n^d
  =\EE Y_n^d+o(\delta_n)
  = \Big(r_n^d+c_n(\sm(r_n)^d-r_n^d)+o(\delta_n)\Big)
  = \Big(r_n^d-\frac{r_n^d-\sm(r_n)^d}{r_n-\sm(r_n)}\delta_n+o(\delta_n)\Big).
\end{align*}
The coefficient
$\frac{r_n^d-\sm(r_n)^d}{r_n-\sm(r_n)}$
is positive and bounded away from zero.  Hence, for all large $n$,
\[
  \EE X_n^d<r_n^d,
\]
contradicting the assumption $\EE X_n^d\ge r_n^d$.  This proves
the lemma.
\end{proof}

\begin{proposition}[Quadratic lower bound for graphons]\label{prop:graphon-quadratic-bound}
There exist constants $C>0$ and $\delta_0>0$, uniform for $r\in K$, such that
if $0<\delta<\delta_0$, $e(W)=r-\delta$, and $t(H,W)\ge r^m$, then
\begin{equation}\label{eq:graphon-quadratic-bound}
  I_{\pc(r)}(W)
  \ge
  J_{\pc(r)}(r)+C\delta^2.
\end{equation}
\end{proposition}

\begin{proof}
Let $X=W(U,V)$, where $(U,V)$ is uniformly distributed on $[0,1]^2$.  Then
$\EE X=e(W)=r-\delta$.  By \eqref{eq:generalized-holder} and feasibility,
\[
  \EE X^d=\int W^d\ge t(H,W)^{d/m} \ge r^d.
\]
By \Cref{lem:scalar-quadratic-bound}, we conclude the proof.
\end{proof}

\subsection{A matching quadratic upper bound}

We now prove that the quadratic lower bound from the previous subsection has
the correct order.  For small $\delta>0$, we construct a bipodal graphon with
edge density $r-\delta$, $H$-density $r^m$, and $I_{\pc(r)}$-value only
$O(\delta^2)$ above that of the constant graphon.  The construction starts
from the constant graphon $r$, makes a small block, uses the second contact density
$s=\sm(r)$ on the cross edges, and then adjusts the large-block density to
restore the two constraints.

\begin{lemma}[Quadratic upper bound]\label{lem:bipodal-quadratic-bound}
There exist constants $C'>0$ and $\delta_0>0$ such that
for every $r\in K$ and every $0<\delta<\delta_0$, there is a bipodal graphon
$V_\delta$ satisfying $e(V_\delta)=r-\delta$, $t(H,V_\delta)=r^m$,
and
\begin{equation}\label{eq:bipodal-quadratic-bound}
  I_{\pc(r)}(V_\delta)
  \le
  J_{\pc(r)}(r)+C'\delta^2.
\end{equation}
\end{lemma}

\begin{proof}
Fix $a\in(0,1)$.
For $r\in K$, put $s=\sm(r)$.  Let $V_{r,c,q}$ be the bipodal graphon with
blocks $A$ and $A^c$, where $|A|=c$, and with edge densities $a$ on
$A\times A$, $s$ on $A\times A^c\cup A^c\times A$, and $q$ on
$A^c\times A^c$.  The block $A\times A$ has measure $c^2$, so its contributions
to the edge density and to $I_{\pc(r)}$ are $O(c^2)$.  Also, any term in the
$H$-density involving this block assigns at least two vertices of $H$ to $A$,
and is therefore $O(c^2)$.  Thus, the parameter $a$ does not enter the
linearization of the two constraints with respect to $(c,q)$ at $(c,q)=(0,r)$.

We first solve the constraints uniformly in $r$.  Define
\[
  F(r,\delta,c,q)
  =
  \bigl(e(V_{r,c,q})-(r-\delta),\,
        t(H,V_{r,c,q})-r^m\bigr).
\]
Although $c$ represents a block size only when $c\in[0,1]$, the formulas for
$e(V_{r,c,q})$ and $t(H,V_{r,c,q})$ are polynomial in the block parameters
and edge probabilities.  We use these polynomial formulas to define $F$ also
for real $c$ near $0$.  Since $K\subset(0,1)\setminus\{r_*\}$ and $\sm$ is
analytic on $(0,1)\setminus\{r_*\}$, this gives an analytic extension of $F$
to a neighbourhood of
\[
  \{(r,0,0,r):r\in K\}.
\]
Also,
\[
  F(r,0,0,r)=(0,0)
  \qquad (r\in K).
\]
At $(c,q)=(0,r)$, the first derivatives of the edge density and
$H$-density with respect to $(c,q)$ are as follows.  The edge density is
\[
  e(V_{r,c,q})=c^2a+2c(1-c)s+(1-c)^2q,
\]
and hence
\[
  \frac{\partial e(V_{r,c,q})}{\partial c}\bigg|_{(0,r)}=2(s-r),
  \qquad
  \frac{\partial e(V_{r,c,q})}{\partial q}\bigg|_{(0,r)}=1,
\]
where the notation $|_{(0,r)}$ means evaluation at $(c,q)=(0,r)$.

For the $H$-density derivative in $c$, write $n=|V(H)|=2m/d$ and expand the
homomorphism density according to how many vertices of $H$ are assigned to the
small block $A$.  If no vertex is assigned to $A$, the contribution is
$(1-c)^nq^m$.  If exactly one vertex is assigned to $A$, then, since every
vertex of $H$ has degree $d$, the $d$ incident edges have density $s$ and the
remaining $m-d$ edges have density $q$; summing over the $n$ choices of this
vertex gives
\[
  n c(1-c)^{n-1}s^d q^{m-d}.
\]
Assignments with at least two vertices in $A$ have total weight $O(c^2)$, so
they do not affect the derivative with respect to $c$ at $c=0$.  Therefore,
\[
  t(H,V_{r,c,q})
  =
  (1-c)^nq^m
  +n c(1-c)^{n-1}s^d q^{m-d}
  +O(c^2),
\]
and evaluating the derivative at $(c,q)=(0,r)$ gives
\[
  \frac{\partial t(H,V_{r,c,q})}{\partial c}\bigg|_{(0,r)}
  =
  -nr^m+ns^d r^{m-d}
  =
  \frac{2m}{d}r^{m-d}(s^d-r^d).
\]
Also, at $c=0$, the graphon is the constant graphon with value $q$, so
$t(H,V_{r,0,q})=q^m$ and hence
\[
  \frac{\partial t(H,V_{r,c,q})}{\partial q}\bigg|_{(0,r)}=m r^{m-1}.
\]
Thus, the Jacobian of $F$ with respect to $(c,q)$ at $(r,0,0,r)$ is
\[
  \begin{pmatrix}
    2(s-r) & 1 \\
    \frac{2m}{d}r^{m-d}(s^d-r^d) & m r^{m-1}
  \end{pmatrix}.
\]
Set
\[
  D_d(r,s)=\frac{s^d-r^d}{d r^{d-1}}-(s-r).
\]
Since $u\mapsto u^d$ is strictly convex and $s\ne r$, we have $D_d(r,s)>0$.
Then, the determinant of the Jacobian is
\[
  -\frac{2m}{d}r^{m-d}(s^d-r^d)+2(s-r)m r^{m-1}
  =-2m r^{m-1}D_d(r,s)<0.
\]
Because $K$ is compact, $r$ stays in a compact subset of $(0,1)$, $s=\sm(r)$
stays in a compact subset of $(0,1)$, and $|s-r|$ is uniformly bounded away
from zero on $K$.  Since $D_d$ is continuous and positive on the compact set
\[
  \{(r,\sm(r)):r\in K\},
\]
the quantity $D_d(r,\sm(r))$ is also uniformly bounded away from zero on $K$.
The above determinant is therefore uniformly bounded away from zero.

By the analytic implicit function theorem, applied with $(r,\delta)$ as
parameters and $(c,q)$ as unknowns, the constraints can be solved locally by an
analytic map of the parameters.  More explicitly, for each $\bar r\in K$, we get an
interval $I_{\bar r}$ containing $\bar r$, a number $\eta_{\bar r}>0$, an open
neighbourhood $W_{\bar r}$ of $(0,\bar r)$ in the $(c,q)$-plane, and an analytic map
\[
  \Xi_{\bar r}(r,\delta)
  =(c_{\bar r}(r,\delta),q_{\bar r}(r,\delta))
\]
defined for $(r,\delta)\in I_{\bar r}\times(-\eta_{\bar r},\eta_{\bar r})$, such that
$F(r,\delta,\Xi_{\bar r}(r,\delta))=(0,0)$ and this is the unique solution with
$(c,q)\in W_{\bar r}$.  After taking $I_{\bar r}$ smaller if necessary, we may assume
that $(0,r)\in W_{\bar r}$ for every $r\in I_{\bar r}$.  Since
$F(r,0,0,r)=(0,0)$, the local uniqueness in $W_{\bar r}$ then gives
\[
  \Xi_{\bar r}(r,0)=(0,r)\qquad (r\in I_{\bar r}).
\]

We now explain why these local solution maps agree on overlaps after choosing a
uniformly small range of $\delta$.  For each $\bar r\in K$, choose an open
interval $J_{\bar r}$ containing $\bar r$ such that
$\overline{J_{\bar r}}\subset I_{\bar r}$.
By compactness, choose $r_1,\ldots,r_k\in K$ such that the intervals
$J_i:=J_{r_i}$ cover $K$.  Write
\[
  I_i=I_{r_i},\qquad
  \eta_i=\eta_{r_i},\qquad
  W_i=W_{r_i},\qquad
  \Xi_i=\Xi_{r_i}.
\]
We claim that, after choosing $\delta_0>0$ sufficiently small,
\[
  \Xi_i(r,\delta)\in W_j
  \qquad
  \text{for all }i,j
  \text{ and all }(r,\delta)\in
  (J_i\cap J_j)\times(-\delta_0,\delta_0).
\]
To see this, fix an ordered pair $(i,j)$ with $J_i\cap J_j\ne\emptyset$.  Since
$\overline{J_j}\subset I_j$ and $(0,r)\in W_j$ for every $r\in I_j$, the
compact set
\[
  \Gamma_{ij}
  =
  \{(0,r):r\in \overline{J_i}\cap\overline{J_j}\}
\]
is contained in the open set $W_j$.  Hence, a small fixed neighbourhood of
$\Gamma_{ij}$ is still contained in $W_j$.  For every
$r\in\overline{J_i}\cap\overline{J_j}$, we have $r\in I_i$, and hence
$\Xi_i(r,0)=(0,r)$.  By continuity of $\Xi_i$, uniformly for
$r\in \overline{J_i}\cap\overline{J_j}$, the values $\Xi_i(r,\delta)$ remain
in this neighbourhood when $|\delta|$ is sufficiently small.  This proves the
claim for the ordered pair $(i,j)$; taking the minimum over the finitely many
ordered pairs, and also requiring $\delta_0<\min_i\eta_i$, proves the claim
simultaneously for all overlaps.

Now let $r\in J_i\cap J_j$ and $|\delta|<\delta_0$.  By the claim,
$\Xi_i(r,\delta)\in W_j$, while $\Xi_j(r,\delta)\in W_j$ by construction.
Both maps solve $F=0$, so the local uniqueness in $W_j$ obtained above gives
$\Xi_i(r,\delta)=\Xi_j(r,\delta)$.  Thus, the local solution
maps are compatible on overlaps and patch together to an analytic map
\[
  \Xi(r,\delta)=(c(r,\delta),q(r,\delta))
\]
on a neighbourhood of $K\times\{0\}$.  The patched map satisfies
\[
  c(r,0)=0,\qquad q(r,0)=r,
\]
and solves the constraints
\begin{equation}\label{eq:comparison-graphon-constraints}
  e(V_{r,c(r,\delta),q(r,\delta)})=r-\delta,
  \qquad
  t(H,V_{r,c(r,\delta),q(r,\delta)})=r^m.
\end{equation}

Differentiating \eqref{eq:comparison-graphon-constraints} with respect to
$\delta$ at $\delta=0$ gives, for every $r\in K$, the linear system
\[
  \partial_\delta q(r,0)+2(s-r)\partial_\delta c(r,0)=-1,
\]
and
\[
  m r^{m-1}\partial_\delta q(r,0)
  +\frac{2m}{d}r^{m-d}(s^d-r^d)\partial_\delta c(r,0)=0.
\]
Solving this system gives
\[
  \partial_\delta c(r,0)=\frac1{2D_d(r,s)},
  \qquad
  \partial_\delta q(r,0)=-\frac{s^d-r^d}{d r^{d-1}D_d(r,s)}
  =-\frac{2(s^d-r^d)}{d r^{d-1}}\partial_\delta c(r,0).
\]
Here the derivatives are with respect to $\delta$ at fixed $r$.  Since
$c(r,\delta)$ and $q(r,\delta)$ are jointly analytic on a neighbourhood of the
compact set $K\times\{0\}$,
Taylor's theorem in the $\delta$ variable has remainders that are uniform in
$r$.  In particular,
\begin{equation}\label{eq:comparison-block-size}
  c(r,\delta)=\frac{\delta}{2D_d(r,s)}+O_K(\delta^2),
\end{equation}
\begin{equation}\label{eq:comparison-block-density}
  q(r,\delta)=
  r-\frac{2(s^d-r^d)}{d r^{d-1}}c(r,\delta)+O_K(c(r,\delta)^2).
\end{equation}
In particular, after shrinking $\delta_0$ uniformly in $r\in K$, we have
\[
  0\le c(r,\delta)\le1,
  \qquad
  0\le q(r,\delta)\le1
  \qquad (0<\delta<\delta_0).
\]
Therefore $V_\delta:=V_{r,c(r,\delta),q(r,\delta)}$ is a valid bipodal graphon
satisfying the two constraints in \eqref{eq:comparison-graphon-constraints}.

It remains to estimate the $I_{p_0}$-value, again uniformly in $r$.  Let
$p_0=\pc(r)$.  We have
\[
  I_{p_0}(V_{r,c,q})
  =
  c^2J_{p_0}(a)+2c(1-c)J_{p_0}(s)+(1-c)^2J_{p_0}(q).
\]
By \eqref{eq:comparison-block-density}, we have $q-r=O_K(c)$.  We now expand the last
expression to first order in the two small quantities $c$ and $q-r$.  The
quantities $a$, $r$, $s=\sm(r)$, and
$p_0=\pc(r)$ already range over compact subsets of $(0,1)$ as $r$ ranges over
$K$.  After shrinking $\delta_0$ uniformly, $q=q(r,\delta)$ also stays in a
compact subset of $(0,1)$.  Thus, $r$, $q$, and $p_0$ are bounded away from
both $0$ and $1$, uniformly for $r\in K$.  This gives the uniform bounds on the
derivatives of $J_{p_0}$ that are needed for the $O_K$ remainder in Taylor's
theorem.  Taylor's theorem therefore gives, uniformly in $r\in K$,
\[
  J_{p_0}(q)
  =
  J_{p_0}(r)+J_{p_0}'(r)(q-r)+O_K((q-r)^2)
  =
  J_{p_0}(r)+J_{p_0}'(r)(q-r)+O_K(c^2).
\]
Also, $J_{p_0}(s)$ is uniformly bounded for $r\in K$, so
\[
  2c(1-c)J_{p_0}(s)
  =
  2cJ_{p_0}(s)+O_K(c^2),
\]
and
\[
  (1-c)^2J_{p_0}(q)
  =
  (1-2c+O(c^2))
  \bigl(J_{p_0}(r)+J_{p_0}'(r)(q-r)+O_K(c^2)\bigr).
\]
Because $q-r=O_K(c)$, the product of $-2c$ with
$J_{p_0}'(r)(q-r)$ is already $O_K(c^2)$.  Hence
\[
  (1-c)^2J_{p_0}(q)
  =
  J_{p_0}(r)+J_{p_0}'(r)(q-r)-2cJ_{p_0}(r)+O_K(c^2).
\]
The remaining term $c^2J_{p_0}(a)$ is $O_K(c^2)$.  Combining these estimates
and subtracting $J_{p_0}(r)$ gives
\[
  I_{p_0}(V_{r,c,q})-J_{p_0}(r)
  =
  J_{p_0}'(r)(q-r)+2c(J_{p_0}(s)-J_{p_0}(r))+O_K(c^2).
\]
Using \eqref{eq:comparison-block-density}, the first-order term equals
\[
  c\left[
    J_{p_0}'(r)\left(-\frac{2(s^d-r^d)}{d r^{d-1}}\right)
    +2(J_{p_0}(s)-J_{p_0}(r))
  \right].
\]
Since the supporting line $\ell_r$ touches the graph of $\varphi_{p_0,d}$ at
both contact points (conditions (M3) and (M4) of
\Cref{thm:scalar-lz-boundary}), we have
\[
  J_{p_0}(s)-J_{p_0}(r)
  =\frac{J_{p_0}'(r)}{d r^{d-1}}(s^d-r^d).
\]
Thus, the first-order term cancels.  Since the edge probabilities
$p_0,a,s,q,r$ remain in a compact subset of $(0,1)$ for $r\in K$ and
$0<\delta<\delta_0$, and since $D_d(r,\sm(r))$ is bounded away from zero on
$K$, the expansion \eqref{eq:comparison-block-size} gives $c=O_K(\delta)$.  Hence, the
remaining terms are
$O_K(c^2)=O_K(\delta^2)$.  Increasing the constant if necessary, we obtain
\[
  I_{\pc(r)}(V_\delta)
  \le
  J_{\pc(r)}(r)+C'\delta^2
\]
with a single constant $C'$ valid for all $r\in K$ and
$0<\delta<\delta_0$.  This proves \eqref{eq:bipodal-quadratic-bound}.
\end{proof}

Recall that the boundary excess \eqref{eq:boundary-excess} is defined only for
$\delta>0$, because the Kenyon--Radin--Ren--Sadun regime of \Cref{thm:krrs-analytic-extension}
requires a strictly positive $H$-density surplus $\tau-\eps^m$.  The theorem
below therefore produces, in addition to the quadratic expansion, the
real-analytic \emph{extension} of $G_r$ to a two-sided neighbourhood of
$\delta=0$; it is that extension whose $\delta$-derivatives appear in
\Cref{sec:local-optimizer-structure}.

\begin{theorem}[Universal positivity of the scalar second variation]\label{thm:positive-second-variation}
There exist $\overline\delta>0$, an open neighbourhood $\mathcal N$ of
$K\times\{0\}$ in $\RR^2$, and a real-analytic function
$G:\mathcal N\to\RR$ such that
\begin{equation}\label{eq:boundary-excess-extension}
  G(r,\delta)=G_r(\delta)
  \qquad(r\in K,\ 0<\delta<\overline\delta),
\end{equation}
and such that, setting $A_H(r):=\partial_\delta^2G(r,0)$, the following hold.
\begin{enumerate}[label=(\roman*)]
  \item $G(r,0)=0$ and $\partial_\delta G(r,0)=0$ for every $r\in K$.
  \item $A_H$ is real-analytic and $A_H(r)>0$ for every $r\in K$; in
  particular $a_0:=\min_{r\in K}A_H(r)>0$.
  \item There is $C_2>0$, depending only on $K$, such that
  \begin{equation}\label{eq:boundary-excess-curvature}
    \bigl|\partial_\delta^2G(r,\delta)-A_H(r)\bigr|\le C_2|\delta|
    \qquad(r\in K,\ |\delta|<\overline\delta).
  \end{equation}
  \item Consequently the boundary excess has the expansion
  \begin{equation}\label{eq:boundary-excess-expansion}
    G_r(\delta)=\frac12A_H(r)\delta^2+O_K(\delta^3)
    \qquad(r\in K,\ 0<\delta<\overline\delta).
  \end{equation}
\end{enumerate}
\end{theorem}

\begin{proof}
By \Cref{thm:krrs-analytic-extension}, the bipodal parameters of the fixed-$(e,t(H,\cdot))$
entropy maximizer appearing in $\mathcal F_{\pc(r),r}(r-\delta)$ depend
analytically on $(\eps,\vartheta)=(r-\delta,\,r^m-(r-\delta)^m)$, and they
extend analytically across the degenerate boundary $\vartheta=0$.  A priori
this is local in $\eps$; but on the finite subcover of $K$ constructed in the
proof of \Cref{lem:fixed-density-bipodality} the local parameter maps agree
on the open set where $\vartheta>0$: there the maximizer is unique up to
relabelling, and the normalization putting its small block first is canonical.
The real-analytic identity theorem then shows that the local extensions also
agree at $\vartheta=0$ and for small negative $\vartheta$.  They therefore
patch to a single analytic map on a neighbourhood of $K\times\{0\}$.  Together with the
analyticity of $\pc$, this lets us extend the boundary excess to a jointly
analytic function, defined on an open neighbourhood $\mathcal N$ of
$K\times\{0\}$,
\[
  G(r,\delta):=\mathcal F_{\pc(r),r}(r-\delta)-J_{\pc(r)}(r),
\]
where for $\delta\le0$ the right-hand side is read through the analytically
extended parameters.  After shrinking $\mathcal N$, choose
$\overline\delta>0$ sufficiently small that
$K\times(-\overline\delta,\overline\delta)\subseteq\mathcal N$ and that the
analytic extension agrees with the boundary excess for $r\in K$ and
$0<\delta<\overline\delta$.  At $\delta=0$, the extended parameters have
$c=0$ and $q_{22}=r$, so they represent the constant graphon $W\equiv r$;
hence $G(r,0)=0$.  This gives \eqref{eq:boundary-excess-extension} and the first equality
in~\textup{(i)}.
Decreasing $\overline\delta$ further if necessary, we may also assume that it
is at most the two radii supplied by \Cref{prop:graphon-quadratic-bound} and
\Cref{lem:bipodal-quadratic-bound}.

For $r\in K$ and $0<\delta<\overline\delta$, let
$\widehat W_{r-\delta,r}$ be the fixed-edge/fixed-$H$ entropy maximizer
appearing in \(\mathcal F_{\pc(r),r}(r-\delta)\).  By
\eqref{eq:reduced-objective-attainment},
\[
  G_r(\delta)
  =
  I_{\pc(r)}(\widehat W_{r-\delta,r})-J_{\pc(r)}(r).
\]
This graphon has edge density $r-\delta$ and $H$-density $r^m$, so
\Cref{prop:graphon-quadratic-bound} gives
\[
  G_r(\delta)\ge C\delta^2.
\]
On the other hand, let $V_\delta$ be the test graphon from
\Cref{lem:bipodal-quadratic-bound}.  It has the same edge density and $H$-density as
$\widehat W_{r-\delta,r}$.  Since $\widehat W_{r-\delta,r}$ minimizes
$I_{\pc(r)}$ on this fixed-constraint slice,
\[
  I_{\pc(r)}(\widehat W_{r-\delta,r})
  \le I_{\pc(r)}(V_\delta)
  \le J_{\pc(r)}(r)+C'\delta^2.
\]
Therefore
\[
  C\delta^2\le G_r(\delta)\le C'\delta^2
  \qquad(r\in K,\ 0<\delta<\overline\delta).
\]

We now read off the remaining assertions.  Since
$\delta\mapsto G(r,\delta)$ is analytic and $G(r,0)=0$, the upper bound
above forces
$\partial_\delta G(r,0)=0$.  To see this,
$|G(r,\delta)|\le C'\delta^2$ for $\delta>0$, so
$\partial_\delta G(r,0)=\lim_{\delta\downarrow0}G(r,\delta)/\delta=0$.  This
is (i).  Consequently the Taylor expansion of $G(r,\cdot)$ at $0$ begins in
degree two, with second coefficient $A_H(r)=\partial_\delta^2G(r,0)$, and
\[
  A_H(r)=\lim_{\delta\downarrow0}\frac{2G_r(\delta)}{\delta^2}\in[2C,2C'].
\]
In particular $A_H(r)\ge2C>0$ uniformly on $K$, and $A_H$ is analytic because
it is the second $\delta$-derivative at $\delta=0$ of the jointly analytic
$G$; this is (ii).

For (iii), shrink $\overline\delta$ once more so that
$K\times[-\overline\delta,\overline\delta]\subseteq\mathcal N$, and let $C_2$ be
$1$ plus the supremum of the continuous function $|\partial_\delta^3G|$ on
that compact set.  In particular, $C_2\ge1$.  The mean value theorem applied to
$\delta\mapsto\partial_\delta^2G(r,\delta)$ on the interval between $0$ and
$\delta$ then gives \eqref{eq:boundary-excess-curvature}.  Finally, Taylor's theorem
with integral remainder together with (i) gives
\[
  G(r,\delta)
  =\int_0^\delta(\delta-t)\,\partial_\delta^2G(r,t)\dd t
  =\frac12A_H(r)\delta^2
   +\int_0^\delta(\delta-t)\bigl(\partial_\delta^2G(r,t)-A_H(r)\bigr)\dd t,
\]
and the last integral is at most $C_2|\delta|^3/6$ in absolute value by
\eqref{eq:boundary-excess-curvature}.  This is (iv), with a cubic constant uniform in
$r\in K$ because $C_2$ is, after using \eqref{eq:boundary-excess-extension} for
$0<\delta<\overline\delta$.
\end{proof}

Although $K$ was fixed in the preceding theorem, the second-variation
function
\[
  r\longmapsto \partial_\delta^2G(r,0)=A_H(r)
\]
constructed there is independent of both $K$ and the chosen analytic
extension.  To see this, consider two such extensions $G_1$ and $G_2$ defined
near the same boundary point $(r_0,0)$.  On a sufficiently small ball contained in
both domains, they agree on the nonempty open subset where $\delta>0$, since
there both equal the boundary excess \eqref{eq:boundary-excess}.  The real-analytic
identity theorem therefore gives $G_1=G_2$ throughout the subset, and hence
\[
  \partial_\delta^2G_1(r_0,0)
  =
  \partial_\delta^2G_2(r_0,0).
\]
Thus, the locally defined second-variation functions patch together to give a
well-defined positive real-analytic function
\begin{equation}\label{eq:second-variation-domain}
  A_H:(0,1)\setminus\{r_*\}\longrightarrow(0,\infty).
\end{equation}